\documentclass[11pt,a4paper]{article}
\usepackage[margin=2.35cm]{geometry}
\usepackage{amsmath,booktabs,graphicx,xcolor,hyperref,array,enumitem}
\graphicspath{{outputs/figures/}{../outputs/figures/}}
\hypersetup{colorlinks=true,linkcolor=blue!55!black,citecolor=blue!55!black,
            urlcolor=blue!55!black}
\setlist{nosep}
\newcommand{\method}[1]{\textsf{#1}}
\newcommand{\NA}{--}

\title{Comparative Unlabelled Anomaly Detection in Daily Climate Time Series}
\author{}
\date{}

\begin{document}
\maketitle

\begin{abstract}
We compare a seasonal burst detector, three discord/profile methods, Isolation
Forest, and three deep reconstruction models on daily temperature,
total precipitation, and pressure. All models are fitted on 1970--1999 and evaluated
on 2000--2023. Score detectors are compared at a detector-specific empirical
99th percentile without rescaling their scores. Raw exceedances are used for
pairwise Jaccard agreement; temporally declustered peaks are used to define and
interpret events. The results show strong dependence on both variable and
anomaly definition. Total precipitation produces the strongest multi-method
agreement, while the z-normalized matrix profile is largely complementary to
amplitude and reconstruction methods. Very small reconstruction loss, most
notably for DLinear, does not imply stronger event validity. In the absence of
labels, the conclusions concern detector behaviour and reproducibility rather
than physical correctness.
\end{abstract}

\section{Objectives}
The study addresses three questions: (i) how detector scores cluster into
events after a common tail-selection rule; (ii) which raw anomalies and events
are shared across detector families; and (iii) how those patterns differ among
temperature, total precipitation, and pressure. The analysis is deliberately
unlabelled. Agreement is therefore interpreted as robustness to method choice,
not as proof that an event is a true physical anomaly.

\section{Experimental protocol}

\subsection{Data preparation and split}
The inputs are daily arrays named \path{temperature_time_series.npy},
\path{preci_time_series.npy}, and \path{pressure_time_series.npy}. A Gregorian
date index is constructed, 29 February is removed when present, and only
complete years in the common period 1 January 1970--31 December 2023 are kept.
This produces 54 years of 365 observations. The first 30 years (1970--1999;
$10{,}950$ days) form the reference period used for fitting, and the remaining
24 years (2000--2023; $8{,}760$ days) form the evaluation period.
Reference-period statistics alone are used for standardization and
reference-neighbour searches.

This split prevents direct evaluation-to-reference leakage. It does not remove the effect
of climatic trends or other distribution shifts, which are part of what the
detectors may score as abnormal. The source metadata and physical units are not
encoded in the current numerical outputs, so physical-unit reconstruction
errors are reported within each variable but not compared across variables.
For readability, physical-data figures convert temperature from K to
$^\circ$C, mean sea-level pressure from Pa to hPa, and total precipitation
from m to mm. These display conversions do not alter detector scores or event
selection.

\subsection{Detectors and fitting}
All score outputs retain their native scale; in particular, they are not
rescaled to $[0,1]$. Larger scores always denote greater abnormality.

\begin{description}[style=nextline]
\item[Seasonal bursts.]
Evaluation aggregates of length 1--15 days are compared with reference aggregates at
the same seasonal position. Absolute standardized departures of at least 3 are
retained and overlapping same-sign candidates are merged. Bursts already
return intervals and do not undergo percentile thresholding.

\item[Subsequence and ensemble detectors.]
The classical subsequence length is $m=15$ days. IMP uses an interval length of
30 days. MP uses z-normalized subsequences and is consequently sensitive to
shape discord; AAMP uses unnormalized distance and retains level and amplitude
information. The nearest reference subsequence is stored for each profile
score. These methods build on the matrix-profile framework
\cite{matrixprofile}. Isolation Forest \cite{iforest} uses 300 trees and seed 0
on overlapping 15-day windows; its output sign is reversed so that large
values are anomalous.

\item[Deep reconstruction detectors.]
AutoEncoder \cite{autoencoder}, DLinear \cite{dlinear}, and TimesNet
\cite{timesnet} receive standardized sequences of length 96. Their overlapping
reconstructions are averaged into a single time series,
and the daily anomaly score is the corresponding overlapping mean squared
error. AutoEncoder hidden-width ratios are 0.5 and 0.25. DLinear uses a
moving-average ratio of 0.1 (an odd kernel of 9 for length 96). TimesNet uses
$d_{\rm model}=16$, $d_{\rm ff}=32$, two layers, top-$k=3$, six convolution
kernels, and dropout 0.1. Each model is trained for 50 epochs using Adam, a
learning rate of $10^{-2}$, batch size 1024, seed 0, and the last 10\% of the
reference-period sequences as a temporal validation set. No early stopping or
validation-based checkpoint selection is used; reported losses are from the
final model.
\end{description}

\subsection{Tail selection and declustering}
For detector $d$, let $s_{d,t}$ denote its evaluation score and let
$q_d=Q_{0.99}(s_d)$ be its empirical 99th percentile. A day is a raw anomaly
when
\begin{equation}
    I_{d,t}=\mathbf{1}\{s_{d,t}>q_d\}.
\end{equation}
Thus the rule selects scores \emph{strictly above} the percentile, rather than
sorting and forcing an exact top-$k$ quota. Ties at the threshold can therefore
make the number differ from one percent. In this run all continuous score
series yield 88 raw exceedances, approximately 1\% of $8{,}760$ evaluation days.

Raw exceedances separated by at most $m=15$ days belong to the same run. Each
run is represented by its maximum-score day. If $N_d$ is the raw count and
$K_d$ the declustered count, we report
\begin{equation}
    \theta_d=K_d/N_d.
\end{equation}
Small $\theta_d$ denotes persistent or recurrent threshold exceedances; large
$\theta_d$ denotes more isolated exceedances. No seasonal threshold is fitted
and no separate set of seasonal peaks is created.

Seasonality is characterized on the complete score series with an STL
decomposition \cite{stl} of period 365, $s_t=T_t+S_t+R_t$. The reported seasonal strength
is
\begin{equation}
    F_S=\max\!\left(0,\min\!\left(1,
       1-\frac{\operatorname{Var}(R_t)}
               {\operatorname{Var}(S_t+R_t)}\right)\right).
\end{equation}
It measures annual organization in the score process; it does not alter the
threshold or the selected anomalies.

\subsection{Agreement and consensus}
Pairwise Jaccard agreement is computed on the \emph{non-declustered} raw flags,
\begin{equation}
 J_{ij}=\frac{|E_i\cap E_j|}{|E_i\cup E_j|},\qquad
 E_i=\{t:I_{i,t}=1\}.
\end{equation}
When two score arrays have different lengths, only their shared dated support
is used; missing trailing dates are not counted as normal observations.

For the event timeline, each declustered representative marks the following 15
days. A selected day for method $i$ is a consensus day if at least one other
method is also active, and the method-level consensus rate is the fraction of
its selected days that are consensus days. The separate consensus-event table
uses a symmetric $\pm15$-day neighbourhood and requires at least three active
score detectors. Its \texttt{n\_detectors} field is the simultaneous count on
the event peak; the detector list records any method participating anywhere in
the merged interval and can therefore be longer.

\subsection{Interpretation views}
Interpretation is performed on declustered representatives. Point detectors
mark the cluster peak and any other raw exceedance from that cluster that falls
inside the displayed window. Deep-model companion points outside the window
are not shown. For IMP, MP, and AAMP, an anomaly is a subsequence: the complete
15-day discord beginning at the cluster peak is displayed in three panels---a
folded evaluation view, the nearest-neighbour folded reference view, and their direct
overlay. For MP, the two folded views retain the original data scale and share
vertical limits; only the direct discord--neighbour comparison is separately
z-normalized. IMP and AAMP retain native scale in all three panels. Bursts
remain intervals. At most five events per method and variable are interpreted.

\section{Results}

\subsection{Peak extraction and score seasonality}
Table~\ref{tab:peaks} reports the complete peak diagnostics. Because the raw
tail budget is the same, differences in event counts are entirely due to the
temporal arrangement of each detector's large scores.

\begin{table}[ht]
\centering
\caption{Raw exceedances, declustered events, event ratio $\theta$, and STL
seasonal strength $F_S$.}
\label{tab:peaks}
\small
\begin{tabular}{llrrrr}
\toprule
Series & Method & Raw & Events & $\theta$ & $F_S$\\
\midrule
Temperature & IMP         & 88 & 17 & .193 & .264\\
            & MP          & 88 & 36 & .409 & .262\\
            & AAMP        & 88 & 17 & .193 & .260\\
            & Isolation Forest & 88 & 13 & .148 & .404\\
            & AutoEncoder & 88 & 67 & .761 & .080\\
            & DLinear     & 88 &  5 & .057 & .803\\
            & TimesNet    & 88 & 15 & .170 & .000\\
\midrule
Total precipitation & IMP         & 88 &  7 & .080 & .189\\
              & MP          & 88 & 28 & .318 & .313\\
              & AAMP        & 88 &  9 & .102 & .160\\
              & Isolation Forest & 88 & 12 & .136 & .177\\
              & AutoEncoder & 88 & 74 & .841 & .339\\
              & DLinear     & 88 & 62 & .705 & .368\\
              & TimesNet    & 88 & 62 & .705 & .303\\
\midrule
Pressure & IMP         & 88 & 18 & .205 & .517\\
         & MP          & 88 & 32 & .364 & .262\\
         & AAMP        & 88 & 14 & .159 & .466\\
         & Isolation Forest & 88 & 19 & .216 & .531\\
         & AutoEncoder & 88 & 44 & .500 & .092\\
         & DLinear     & 88 & 33 & .375 & .199\\
         & TimesNet    & 88 & 52 & .591 & .529\\
\bottomrule
\end{tabular}
\end{table}

Clustering depends strongly on both method and variable. Temperature DLinear
has only five representatives ($0.06$), so its top errors occur in a few
long or recurrent runs; its score is also strongly annual ($F_S=0.80$). In
contrast, temperature AutoEncoder produces 67 relatively isolated events and
weak score seasonality. Total precipitation shows a detector-family split: IMP,
AAMP, and Isolation Forest compress 88 days into only 7--12 events, whereas the
deep models retain 62--74 events. Pressure is intermediate. Its IMP, AAMP,
Isolation Forest, and TimesNet scores all have substantial annual organization
($F_S=.466$--$.531$), which cautions against interpreting a global-tail event
list as seasonally uniform.

\subsection{Agreement}
Table~\ref{tab:agreement} gives the number of events and the percentage of each
method's selected timeline days overlapping at least one other method.

\begin{table}[ht]
\centering
\caption{Declustered event count and method-level consensus rate $C$ (percent).}
\label{tab:agreement}
\small
\begin{tabular}{lrr@{\qquad}rr@{\qquad}rr}
\toprule
& \multicolumn{2}{c}{Temperature} & \multicolumn{2}{c}{Total precipitation}
& \multicolumn{2}{c}{Pressure}\\
Method & Events & $C$ & Events & $C$ & Events & $C$\\
\midrule
Burst             & 58 & 45.4 & 252 & 55.4 & 57 & 46.7\\
IMP               & 17 & 92.5 &   7 &100.0 & 18 & 81.5\\
MP                & 36 & 18.5 &  28 & 43.6 & 32 & 11.9\\
AAMP              & 17 & 85.9 &   9 & 99.3 & 14 & 93.3\\
Isolation Forest  & 13 & 73.3 &  12 & 98.9 & 19 & 76.8\\
AutoEncoder       & 67 & 22.9 &  74 & 82.5 & 44 & 70.5\\
DLinear           &  5 & 57.3 &  62 & 92.0 & 33 & 49.5\\
TimesNet          & 15 & 53.3 &  62 & 77.2 & 52 & 55.3\\
\bottomrule
\end{tabular}
\end{table}

The non-declustered Jaccard matrices make the family structure clearer.
For total precipitation, IMP--AAMP has $J=.743$, AAMP--Isolation Forest $J=.645$,
and IMP--Isolation Forest $J=.530$. The deep family forms a second block:
AutoEncoder--DLinear $J=.544$, AutoEncoder--TimesNet $J=.375$, and
DLinear--TimesNet $J=.354$. Cross-family overlaps are much smaller. For
temperature the largest overlap is IMP--AAMP ($J=.257$), followed by
IMP--Isolation Forest ($J=.222$); for pressure it is again IMP--AAMP
($J=.419$). MP has zero raw overlap with every other score detector for
pressure and nearly zero overlap for temperature. This is consistent with MP
detecting z-normalized shape discord rather than level or amplitude extremes.

Total precipitation has the strongest event-level reproducibility: the mean
method-level consensus rate across the seven score detectors is 84.8\%, versus
62.7\% for pressure and 57.7\% for temperature. Nevertheless, high consensus
rates for IMP and AAMP partly reflect their small, persistent event sets; a high
rate should not be read independently of event count and raw Jaccard.

\begin{figure}[ht]
\centering
\includegraphics[width=.94\textwidth]{cross_series/method_comparison.png}
\caption{Cross-variable comparison of event count, consensus rate, raw
Jaccard agreement, and score seasonality.}
\label{fig:cross}
\end{figure}

The symmetric consensus procedure yields 15 temperature, 41 total-precipitation,
and 36 pressure intervals with at least three simultaneously active score
detectors. Peak concurrence reaches five detectors for temperature and six for
total precipitation and pressure. Examples of maximal concurrence are shown in
Table~\ref{tab:consensus-events}; these dates are candidates for physical
follow-up, not externally validated events.

\begin{table}[ht]
\centering
\caption{Illustrative maximal-concurrence consensus peaks. The method list may
include participation away from the peak within the merged interval.}
\label{tab:consensus-events}
\small
\begin{tabular}{llcl}
\toprule
Series & Peak date & Peak count & Participating methods\\
\midrule
Temperature & 2003-08-02 & 5 & AAMP, DLinear, IMP, MP, iForest\\
            & 2019-07-08 & 5 & AAMP, AE, DLinear, IMP, TimesNet, iForest\\
            & 2023-08-08 & 5 & AAMP, DLinear, IMP, TimesNet, iForest\\
Total precipitation & 2015-07-29 & 6 & AAMP, AE, DLinear, IMP, TimesNet, iForest\\
              & 2015-12-26 & 6 & AAMP, AE, DLinear, IMP, TimesNet, iForest\\
              & 2020-04-26 & 6 & AAMP, AE, DLinear, IMP, TimesNet, iForest\\
Pressure & 2000-04-01 & 6 & AAMP, AE, DLinear, IMP, TimesNet, iForest\\
         & 2010-10-26 & 6 & AAMP, AE, DLinear, IMP, MP, TimesNet, iForest\\
         & 2017-11-29 & 6 & AAMP, AE, DLinear, IMP, TimesNet, iForest\\
\bottomrule
\end{tabular}
\end{table}

\begin{figure}[p]
\centering
\includegraphics[width=.96\textwidth]{temperature/comparison/agreement_timeline.png}
\caption{Temperature event timeline and daily method agreement.}
\end{figure}
\begin{figure}[p]
\centering
\includegraphics[width=.96\textwidth]{precipitation/comparison/agreement_timeline.png}
\caption{Total-precipitation event timeline and daily method agreement.}
\end{figure}
\begin{figure}[p]
\centering
\includegraphics[width=.96\textwidth]{pressure/comparison/agreement_timeline.png}
\caption{Pressure event timeline and daily method agreement.}
\end{figure}

\subsection{Deep reconstruction diagnostics}
Tables~\ref{tab:deep-loss} and~\ref{tab:deep-error} report final standardized
MSE and absolute errors after overlapping reconstructions are mapped back to
the original variable scale.
Reference-fit, validation, and evaluation losses are close enough to give no evidence of a
large final-epoch generalization gap, although only one temporal split and one
seed are available. Temperature TimesNet has the clearest validation increase
(reference fit $.00288$, validation $.00691$), while its evaluation loss lies between them.

\begin{table}[ht]
\centering
\caption{Final deep-model mean squared error on standardized data.}
\label{tab:deep-loss}
\small
\begin{tabular}{llrrr}
\toprule
Series & Model & Reference-fit MSE & Validation MSE & Evaluation MSE\\
\midrule
Temperature & AutoEncoder & .05936 & .05698 & .06204\\
 & DLinear & $2.24\!\times\!10^{-4}$ & $1.03\!\times\!10^{-4}$ & $1.19\!\times\!10^{-4}$\\
 & TimesNet & .00288 & .00691 & .00429\\
\midrule
Total precipitation & AutoEncoder & .53821 & .52794 & .52739\\
 & DLinear & $6.94\!\times\!10^{-14}$ & $6.64\!\times\!10^{-14}$ & $6.76\!\times\!10^{-14}$\\
 & TimesNet & .00723 & .00257 & .00361\\
\midrule
Pressure & AutoEncoder & .18981 & .17094 & .20311\\
 & DLinear & $1.19\!\times\!10^{-6}$ & $8.62\!\times\!10^{-7}$ & $9.10\!\times\!10^{-7}$\\
 & TimesNet & .00854 & .00449 & .00475\\
\bottomrule
\end{tabular}
\end{table}

\begin{table}[ht]
\centering
\caption{Absolute errors of the overlapping reconstruction time series. Values
retain each variable's original scale.}
\label{tab:deep-error}
\small
\begin{tabular}{llrrr}
\toprule
Series & Model & Median $|e|$ & Mean $|e|$ & Maximum $|e|$\\
\midrule
Temperature & AutoEncoder & .9151 & 1.1165 & 6.8143\\
 & DLinear & .0030 & .0033 & .0276\\
 & TimesNet & .0420 & .0959 & 2.8480\\
\midrule
Total precipitation & AutoEncoder & .00977 & .02152 & .42584\\
 & DLinear & $1.27\!\times\!10^{-8}$ & $1.87\!\times\!10^{-8}$ & $6.67\!\times\!10^{-7}$\\
 & TimesNet & .000754 & .001264 & .077997\\
\midrule
Pressure & AutoEncoder & 201.17 & 252.06 & 1691.08\\
 & DLinear & .3335 & .4727 & 4.2362\\
 & TimesNet & 20.7288 & 28.6286 & 576.405\\
\bottomrule
\end{tabular}
\end{table}

DLinear reconstructs all three series almost exactly, most strikingly for
total precipitation. This is a property of the learned reconstruction mapping and
the smooth/zero-heavy input, not evidence that DLinear has discovered a
superior anomaly definition. Percentile thresholding still ranks its very
small residuals and selects 88 days. The selected dates can consequently be
sensitive to numerical or subtle modelling differences. More generally,
losses cannot be ranked across architectures as a substitute for anomaly
quality: each architecture has a different reconstruction bias, and the
physical-scale errors cannot be ranked across variables without common units
and scientific tolerances.

\section{Analysis by variable}

\subsection{Temperature}
Temperature has the weakest score-detector consensus of the three variables.
IMP, AAMP, and Isolation Forest form a modest amplitude/local-pattern group,
whereas MP and the deep detectors contribute largely distinct raw days. The
five DLinear events arise from strongly seasonal and temporally clustered
residuals; the 67 AutoEncoder events are far more dispersed. The maximal
multi-method peaks occur in boreal summer in the current record, but no heat
event attribution should be made without the station/grid metadata and an
independent catalogue. The event-month histogram and folded interpretations
should be used to distinguish a seasonal opportunity effect from exceptional
within-season values.

\subsection{Total precipitation}
Total precipitation produces both the most burst intervals (252) and the highest
cross-method event consensus. Its raw Jaccard matrix separates an
amplitude/local-window family (IMP, AAMP, Isolation Forest) from a reconstruction
family (AutoEncoder, DLinear, TimesNet). This suggests that the same broad
episodes can contain two different signals: concentrated local extremes and
distributed reconstruction departures. MP remains more specific, with only
43.6\% timeline consensus and low raw Jaccard. Because total precipitation is often
zero-heavy and intermittent, the nearly exact DLinear reconstruction and the
global percentile threshold require particular caution; inspection should
report wet-day frequency and accumulation magnitude before calling these
events meteorologically exceptional.

\subsection{Pressure}
Pressure has intermediate event consensus but the highest mean score
seasonality. IMP, AAMP, Isolation Forest, and TimesNet each have $F_S$ near or
above .47, so their global 99th percentile emphasizes seasons in which their
score distributions are naturally broader. AAMP has the highest timeline
consensus (93.3\%), while MP is almost entirely method-specific (11.9\% and
zero raw overlap with the other score detectors). The distinction is
scientifically plausible as a distinction between level/amplitude anomalies
and unusual normalized shape, but confirming it requires examining pressure
tendency and synoptic context around the displayed subsequences.

\section{Validity, sensitivity, and next analyses}
The present run supports a comparative detector analysis, not a calibrated
estimate of anomaly-detection accuracy. Its main limitations are:
\begin{itemize}
\item no labelled truth or independent event catalogue;
\item one reference/evaluation boundary, one random seed, and final-epoch deep models;
\item dependence among overlapping windows and among days within a cluster;
\item fixed 15-day classical and 96-day deep contexts;
\item a single global 99th-percentile rule despite variable score seasonality;
\item no uncertainty intervals for event counts, Jaccard, or consensus; and
\item missing source metadata, units, quality-control history, and spatial
support in the report inputs.
\end{itemize}

The priority sensitivity analysis is a grid over percentile
($.975,.99,.995$) and declustering distance (for example 7, 15, and 30 days),
with year-block bootstrap intervals for Jaccard and consensus. Deep models
should be repeated over multiple seeds and compared with a simple persistence
or seasonal-climatology reconstruction baseline. External event validation
criteria should be fixed before inspecting candidate dates. These additions
would reveal whether the observed family structure is stable or is an artifact
of one tail budget and one temporal scale.

\section{Reproducibility}
The numerical results used here are stored below \path{outputs/figures/} in
one directory per variable. The principal machine-readable files are:
\begin{itemize}
\item \path{peak_diagnostics.csv}, \path{detector_jaccard.csv}, and
\path{agreement_summary.csv};
\item \path{consensus_events.csv} and \path{deep_losses.csv}; and
\item \path{deep_reconstruction_summary.csv}.
\end{itemize}
Cross-variable summaries and figures are in \path{cross_series/}. A complete
archive should additionally record the repository revision, environment,
hardware, input checksums, runtime, and full source metadata.

\begin{thebibliography}{9}
\bibitem{iforest}
F.~T. Liu, K.~M. Ting, and Z.-H. Zhou,
Isolation Forest,
\emph{Proceedings of the IEEE International Conference on Data Mining}, 2008.

\bibitem{matrixprofile}
C.-C.~M. Yeh et al.,
Matrix Profile I: All Pairs Similarity Joins for Time Series,
\emph{Proceedings of the IEEE International Conference on Data Mining}, 2016.

\bibitem{autoencoder}
G.~E. Hinton and R.~R. Salakhutdinov,
Reducing the Dimensionality of Data with Neural Networks,
\emph{Science}, 313(5786):504--507, 2006.

\bibitem{dlinear}
A. Zeng et al.,
Are Transformers Effective for Time Series Forecasting?,
\emph{Proceedings of the AAAI Conference on Artificial Intelligence}, 2023.

\bibitem{timesnet}
H. Wu et al.,
TimesNet: Temporal 2D-Variation Modeling for General Time Series Analysis,
\emph{International Conference on Learning Representations}, 2023.

\bibitem{stl}
R.~B. Cleveland, W.~S. Cleveland, J.~E. McRae, and I. Terpenning,
STL: A Seasonal-Trend Decomposition Procedure Based on Loess,
\emph{Journal of Official Statistics}, 6(1):3--73, 1990.
\end{thebibliography}

\end{document}
